\documentclass[11pt,letterpaper]{article}
\usepackage[utf8]{inputenc}

%----- Configuración del estilo del documento------%
\usepackage{epsfig,graphicx}
\usepackage[left=2cm,right=2cm,top=1.8cm,bottom=2.3cm]{geometry}
\usepackage{fancyhdr}
\usepackage{lastpage}
\usepackage{url}
\pagestyle{fancy}
\fancyhf{}
\rfoot{\textcolor{white}{\textit{Página \thepage \hspace{1pt} de \pageref{LastPage}}}}

%------ Paquetes matemáticos básicos --------%
\usepackage{amsmath,amssymb,amsthm}

%------ Idioma --------%
\usepackage[spanish]{babel}

%------ Colores y gráficos --------%
\usepackage{xcolor}
\usepackage[table]{xcolor}
\usepackage{colortbl}
\usepackage{graphicx}
\usepackage{array, multirow, multicol, tabularx}
\usepackage{tcolorbox}
\tcbuselibrary{skins, breakable}

%------ Modo oscuro --------%
\pagecolor{black}
\color{white}

%------ Colores personalizados --------%
\definecolor{B}{HTML}{000000}
\definecolor{G}{HTML}{CCCCCC}
\definecolor{R2}{HTML}{ff6f6f}
\definecolor{A2}{HTML}{7fbfff}
\definecolor{V2}{HTML}{aaffaa}

%------ Colores neón --------%
\definecolor{neongreen}{rgb}{0.22, 1.0, 0.22}
\definecolor{neonpink}{rgb}{1.0, 0.2, 0.6}
\definecolor{neonblue}{rgb}{0.3, 0.9, 1.0}
\definecolor{neonyellow}{rgb}{1.0, 1.0, 0.3}
\definecolor{neonorange}{rgb}{1.0, 0.5, 0.0}
\definecolor{neonpurple}{rgb}{0.8, 0.2, 1.0}

%------ Comandos matemáticos --------%
\newcommand{\R}{\mathbb{R}}
\newcommand{\C}{\mathcal{C}}
\newcommand{\N}{\mathbb{N}}
\newcommand{\Z}{\mathbb{Z}}
\newcommand{\Q}{\mathbb{Q}}
\renewcommand{\theenumi}{\Roman{enumi}}
\renewcommand{\labelenumi}{{\theenumi}.}

%------ Teoremas --------%
\newtheorem{theorem}{Teorema}[section]
\newtheorem{corollary}{Corolario}[theorem]
\newtheorem{lemma}[theorem]{Lema}

%------ TikZ y librerías gráficas --------%
\usepackage{tikz}
\usepackage{tikz-3dplot}
\usetikzlibrary{arrows.meta} % Flechas como Stealth
\usetikzlibrary{graphs}      % Entorno \graph
\usetikzlibrary{quotes}      % Etiquetas en aristas
\usetikzlibrary{positioning} % Posicionamiento relativo
\usetikzlibrary{calc}        % Coordenadas calculadas
\usetikzlibrary{decorations.pathmorphing} % Curvas decoradas
\usetikzlibrary{shapes.geometric} % Formas como poliedros
\tdplotsetmaincoords{70}{30}

%------ Hipervínculos --------%
\usepackage{hyperref}
\hypersetup{
  colorlinks=true,
  linkcolor=neonpink,
  urlcolor=neonblue,
  citecolor=neongreen
}

\usetikzlibrary{calc}

% Fondo negro
\pagecolor{black}

% Colores análogos neón (más saturados)
\definecolor{neonblue}{RGB}{0,150,255}
\definecolor{neoncyan}{RGB}{0,255,255}
\definecolor{neongreen}{RGB}{0,255,128}

% Estilos
\tikzset{
  caraA/.style={draw=neonblue, line width=1.4pt, fill=neonblue!25, opacity=0.65},
  caraB/.style={draw=neoncyan, line width=1.4pt, fill=neoncyan!25, opacity=0.65},
  caraC/.style={draw=neongreen, line width=1.4pt, fill=neongreen!25, opacity=0.65},
  vertice/.style={fill=neoncyan},
  dashedline/.style={draw=neonblue, line width=1pt, dash pattern=on 2pt off 2pt},
  glow/.style={line width=3pt, draw=#1!30, opacity=0.6}
  }
\tikzset{
  cara/.style={draw=neoncyan, line width=2pt, fill=neoncyan!15, opacity=0.4},
  vertice/.style={fill=neoncyan, circle, inner sep=2pt},
  glow/.style={line width=4pt, draw=#1!40, opacity=0.5},
  arista-vertical/.style={draw=neonpink, line width=2.5pt, opacity=0.9},
  arista-horizontal/.style={draw=neonyellow, line width=2.2pt, opacity=0.9},
  borde-cara/.style={draw=neoncyan, line width=1.5pt},
  % Nuevo estilo para aristas del fondo (punteadas)
  arista-fondo/.style={draw=neoncyan!60, line width=1.8pt, opacity=0.7, 
                       dash pattern=on 4pt off 3pt}
}


\begin{document}
% Requiere: \usepackage{tikz}
% Opcional: \usetikzlibrary{arrows.meta}



\begin{tikzpicture}[scale=1.0, every node/.style={font=\small}]
  %----- Estilos -----
  \tikzset{
    vertex/.style={circle, fill=black, draw=black, inner sep=0pt, minimum size=4pt},
    Llabel/.style={anchor=east},
    Rlabel/.style={anchor=west},
    eblue/.style={line width=1pt, draw=blue!70},
    egreen/.style={line width=1pt, draw=green!70!black},
    ered/.style={line width=1pt, draw=red!75!black},
    boundary/.style={rounded corners=6pt, line width=0.9pt, draw=black!60},
    flow/.style={-{Stealth[length=3pt]}, line width=0.8pt, draw=black!60}
  }

  %----- Rectángulo que representa el toro -----
  \draw[boundary] (-0.8,-3.0) rectangle (6.8,3.0);

  % Flechas de identificación (toro)
  \draw[flow] (3.0,-3.0) -- (3.8,-3.0); % abajo derecha
  \draw[flow] (3.8,3.0) -- (3.0,3.0);   % arriba izquierda
  \draw[flow] (-0.8,-1.0) -- (-0.8,-1.8); % izquierda abajo
  \draw[flow] (6.8,1.0) -- (6.8,1.8);     % derecha arriba

  %----- Vértices izquierda -----
  \node[vertex] (L1) at (0,2.0) {};
  \node[vertex] (L2) at (0,0.0) {};
  \node[vertex] (L3) at (0,-2.0) {};
  \node[Llabel] at ($(L1)+(-0.2,0)$) {$L_1$};
  \node[Llabel] at ($(L2)+(-0.2,0)$) {$L_2$};
  \node[Llabel] at ($(L3)+(-0.2,0)$) {$L_3$};

  %----- Vértices derecha -----
  \node[vertex] (R1) at (6,2.0) {};
  \node[vertex] (R2) at (6,0.0) {};
  \node[vertex] (R3) at (6,-2.0) {};
  \node[Rlabel] at ($(R1)+(0.2,0)$) {$R_1$};
  \node[Rlabel] at ($(R2)+(0.2,0)$) {$R_2$};
  \node[Rlabel] at ($(R3)+(0.2,0)$) {$R_3$};

  %----- Aristas azules (directas) -----
  \draw[eblue] (L1) -- (R1);
  \draw[eblue] (L2) -- (R2);
  \draw[eblue] (L3) -- (R3);

  %----- Aristas verdes (diagonales cruzadas) -----
  \draw[egreen] (L1) -- (R2);
  \draw[egreen] (L2) -- (R3);
  \draw[egreen] (L3) -- (R1);

  %----- Aristas rojas (curvas internas) -----
  \draw[ered] (L1) .. controls (3,2.5) .. (R3);
  \draw[ered] (L2) .. controls (3,0.5) .. (R1);
  \draw[ered] (L3) .. controls (3,-2.5) .. (R2);

  %----- Etiqueta central -----
  \node[font=\normalsize] at (3.0, 3.3) {$K_{3,3}$};
\end{tikzpicture}

\begin{center}
\begin{tikzpicture}[thick]

% Nivel superior
\foreach \v/\a in {A/10,B/82,C/154,D/226,E/298}{
  \coordinate (\v) at ({4*cos(\a)},{.8*sin(\a)});
  \fill[vertice] (\v) circle [radius=0.12];
}
\draw[glow=neongreen] (A)--(B)--(C)--(D)--(E)--cycle;
\draw[neongreen, line width=1.4pt, fill=neoncyan!20, opacity=0.7] (A)--(B)--(C)--(D)--(E)--cycle;

% Segundo nivel
\begin{scope}[shift={(0,-3)}]
  \foreach \v/\a in {A'/10,B'/82,C'/154,D'/226,E'/298}{
    \coordinate (\v) at ({6*cos(\a)},{1.2*sin(\a)});
    \fill[vertice] (\v) circle [radius=0.12];
  }
\end{scope}

% Tercer nivel
\begin{scope}[shift={(0,-5.5)}]
  \foreach \v/\a in {F/190,G/262,H/334,I/406,J/478}{
    \coordinate (\v) at ({6*cos(\a)},{1.2*sin(\a)});
    \fill[vertice] (\v) circle [radius=0.12];
  }
\end{scope}

% Cuarto nivel
\begin{scope}[shift={(0,-9)}]
  \foreach \v/\a in {F'/190,G'/262,H'/334,I'/406,J'/478}{
    \coordinate (\v) at ({4*cos(\a)},{.8*sin(\a)});
    \fill[vertice] (\v) circle [radius=0.12];
  }
  \draw[glow=neonblue] (F')--(G')--(H')--(I')--(J')--cycle;
  \draw[dashedline,fill=neonblue!30,opacity=0.6] (F')--(G')--(H')--(I')--(J')--cycle;
\end{scope}

% Caras delineadas con glow
\draw[glow=neonblue] (C)--(D)--(D')--(F)--(C')--cycle;
\draw[caraA] (C)--(D)--(D')--(F)--(C')--cycle;

\draw[glow=neoncyan] (D)--(E)--(E')--(G)--(D')--cycle;
\draw[caraB] (D)--(E)--(E')--(G)--(D')--cycle;

\draw[glow=neongreen] (E)--(A)--(A')--(H)--(E')--cycle;
\draw[caraC] (E)--(A)--(A')--(H)--(E')--cycle;

\draw[glow=neonblue] (F)--(D')--(G)--(G')--(F')--cycle;
\draw[caraA] (F)--(D')--(G)--(G')--(F')--cycle;

\draw[glow=neoncyan] (G)--(E')--(H)--(H')--(G')--cycle;
\draw[caraB] (G)--(E')--(H)--(H')--(G')--cycle;

% Líneas punteadas con glow
\draw[glow=neonblue] (C') -- (J);
\draw[dashedline] (C') -- (J);

\draw[glow=neonblue] (B') -- (B);
\draw[dashedline] (B') -- (B);

\draw[glow=neonblue] (A') -- (I);
\draw[dashedline] (A') -- (I);

\draw[glow=neonblue] (J') -- (J) -- (B') -- (I) -- (I');
\draw[dashedline] (J') -- (J) -- (B') -- (I) -- (I');
\end{tikzpicture}
\end{center}
%cubo
\begin{center}
\begin{tikzpicture}[scale=4]

% Vértices de un tetraedro regular centrado
\coordinate (A) at (0,0,0);
\coordinate (B) at (2,0,0);
\coordinate (C) at (1,1.732,0);
\coordinate (D) at (1,0.577,1.633);

% Rotación ligera para ver la base
\begin{scope}[rotate around={20:(1,0.8)}]

  % Vértices brillantes finos
  \foreach \v in {A,B,C,D}{
    \fill[vertice] (\v) circle (0.04);
  }

  % Caras con glow y traslúcidas
  \draw[glow=neonblue] (A)--(B)--(C)--cycle;
  \draw[caraC] (A)--(B)--(C)--cycle;

  \draw[glow=neoncyan] (A)--(B)--(D)--cycle;
  \draw[caraB] (A)--(B)--(D)--cycle;

  \draw[glow=neongreen] (B)--(C)--(D)--cycle;
  \draw[caraA] (B)--(C)--(D)--cycle;

  \draw[glow=neonblue] (C)--(A)--(D)--cycle;
  \draw[caraC] (C)--(A)--(D)--cycle;

\end{scope}
\end{tikzpicture}
\end{center}
\begin{center}
\begin{tikzpicture}[scale=3.5]

% Vértices del cubo
\coordinate (A) at (0,0,0);
\coordinate (B) at (2,0,0);
\coordinate (C) at (2,2,0);
\coordinate (D) at (0,2,0);
\coordinate (E) at (0,0,2);
\coordinate (F) at (2,0,2);
\coordinate (G) at (2,2,2);
\coordinate (H) at (0,2,2);

% Vértices brillantes más finos
\foreach \v in {A,B,C,D,E,F,G,H}{
  \fill[vertice] (\v) circle (0.04); % mucho más fino
}

% Caras con glow
\draw[glow=neonblue] (A)--(B)--(C)--(D)--cycle;
\draw[caraC] (A)--(B)--(C)--(D)--cycle;

\draw[glow=neoncyan] (E)--(F)--(G)--(H)--cycle;
\draw[caraC] (E)--(F)--(G)--(H)--cycle;

\draw[glow=neongreen] (A)--(B)--(F)--(E)--cycle;
\draw[caraA] (A)--(B)--(F)--(E)--cycle;

\draw[glow=neonblue] (B)--(C)--(G)--(F)--cycle;
\draw[caraC] (B)--(C)--(G)--(F)--cycle;

\draw[glow=neoncyan] (C)--(D)--(H)--(G)--cycle;
\draw[caraB] (C)--(D)--(H)--(G)--cycle;

\draw[glow=neongreen] (D)--(A)--(E)--(H)--cycle;
\draw[caraB] (D)--(A)--(E)--(H)--cycle;

\end{tikzpicture}
\end{center}

\begin{center}
%octaedro
\begin{tikzpicture}[tdplot_main_coords, scale=1.2]
\usetikzlibrary{backgrounds}

% Escala
\def\S{3.0}

% ROTACIÓN MUY SUAVE del cuadrado medio
\pgfmathsetmacro{\rotAng}{5}
\pgfmathsetmacro{\cosRot}{cos(\rotAng)}
\pgfmathsetmacro{\sinRot}{sin(\rotAng)}

% Vértices superiores e inferiores
\coordinate (T) at (0,0,{1.8*\S});  % Top
\coordinate (B) at (0,0,{-1.8*\S}); % Bottom

% Vértices en el plano XY formando un cuadrado AMPLIADO y LEVEMENTE ROTADO
\def\R{1.2*\S}

% Coordenadas originales
\pgfmathsetmacro{\Ax}{\R}
\pgfmathsetmacro{\Ay}{0}
\pgfmathsetmacro{\Bx}{0}
\pgfmathsetmacro{\By}{\R}
\pgfmathsetmacro{\Cx}{-\R}
\pgfmathsetmacro{\Cy}{0}
\pgfmathsetmacro{\Dx}{0}
\pgfmathsetmacro{\Dy}{-\R}

% Aplicar rotación mínima
\coordinate (A) at ({\Ax*\cosRot - \Ay*\sinRot}, {\Ax*\sinRot + \Ay*\cosRot}, 0);
\coordinate (Bp) at ({\Bx*\cosRot - \By*\sinRot}, {\Bx*\sinRot + \By*\cosRot}, 0);
\coordinate (C) at ({\Cx*\cosRot - \Cy*\sinRot}, {\Cx*\sinRot + \Cy*\cosRot}, 0);
\coordinate (D) at ({\Dx*\cosRot - \Dy*\sinRot}, {\Dx*\sinRot + \Dy*\cosRot}, 0);

% PRIMERO: Dibujar las aristas del FONDO con líneas punteadas
% Identificamos qué aristas están en la parte trasera según la perspectiva actual
% Con esta perspectiva (70,30), las aristas traseras son aproximadamente:
% - Parte de las aristas verticales del lado opuesto
% - Algunas aristas del cuadrado medio

% Aristas verticales del fondo (punteadas)
\draw[arista-fondo] (B) -- (C);  % Vertical trasera izquierda
\draw[arista-fondo] (T) -- (C);  % Vertical trasera izquierda superior

% Aristas del cuadrado medio en el fondo (punteadas)
\draw[arista-fondo] (C) -- (D);  % Lado trasero del cuadrado
\draw[arista-fondo] (D) -- (A);  % Lado diagonal trasero del cuadrado

% SEGUNDO: Dibujar las aristas del PRIMER PLANO (sólidas)
% Aristas verticales delanteras (sólidas)
\foreach \v in {A, Bp}{
  \draw[arista-vertical] (T) -- (\v);
  \draw[arista-vertical] (B) -- (\v);
}

% Arista vertical delantera derecha
\draw[arista-vertical] (T) -- (D);
\draw[arista-vertical] (B) -- (D);

% Aristas horizontales del cuadrado medio delanteras (sólidas)
\draw[arista-horizontal] (A) -- (Bp);
\draw[arista-horizontal] (Bp) -- (C);  % Parcialmente visible

% TERCERO: Caras del octaedro (sobre las aristas)
% Caras delanteras (más opacas)
% Cara superior frontal 1
\draw[glow=neonblue] (T)--(A)--(Bp)--cycle;
\draw[caraA, opacity=0.5] (T)--(A)--(Bp)--cycle;
\draw[borde-cara] (T)--(A)--(Bp)--cycle;

% Cara superior frontal 4
\draw[glow=neonblue] (T)--(D)--(A)--cycle;
\draw[caraB, opacity=0.5] (T)--(D)--(A)--cycle;
\draw[borde-cara] (T)--(D)--(A)--cycle;

% Cara inferior frontal 1
\draw[glow=neonblue] (B)--(A)--(Bp)--cycle;
\draw[caraA, opacity=0.5] (B)--(A)--(Bp)--cycle;
\draw[borde-cara] (B)--(A)--(Bp)--cycle;

% Cara inferior frontal 4
\draw[glow=neonblue] (B)--(D)--(A)--cycle;
\draw[cara, opacity=0.5] (B)--(D)--(A)--cycle;
\draw[borde-cara] (B)--(D)--(A)--cycle;

% Caras traseras (menos opacas para que se vean las aristas punteadas)
% Cara superior trasera 2
\draw[glow=neonblue, opacity=0.3] (T)--(Bp)--(C)--cycle;
\draw[caraC, opacity=0.3] (T)--(Bp)--(C)--cycle;
\draw[borde-cara, opacity=0.3] (T)--(Bp)--(C)--cycle;

% Cara superior trasera 3
\draw[glow=neonblue, opacity=0.3] (T)--(C)--(D)--cycle;
\draw[caraC, opacity=0.3] (T)--(C)--(D)--cycle;
\draw[borde-cara, opacity=0.3] (T)--(C)--(D)--cycle;

% Cara inferior trasera 2
\draw[glow=neonblue, opacity=0.3] (B)--(Bp)--(C)--cycle;
\draw[caraC, opacity=0.3] (B)--(Bp)--(C)--cycle;
\draw[borde-cara, opacity=0.3] (B)--(Bp)--(C)--cycle;

% Cara inferior trasera 3
\draw[glow=neonblue, opacity=0.3] (B)--(C)--(D)--cycle;
\draw[caraC, opacity=0.3] (B)--(C)--(D)--cycle;
\draw[borde-cara, opacity=0.3] (B)--(C)--(D)--cycle;

% Vértices - con brillo intenso
\foreach \v in {T,B}{
  \fill[neonpink!90] (\v) circle (2.5pt);
  \fill[white] (\v) circle (1pt);
}

% Vértices del cuadrado medio
\foreach \v in {A,Bp,C,D}{
  \fill[neonyellow!90] (\v) circle (2.5pt);
  \fill[white!80] (\v) circle (0.8pt);
}

% Efectos de resplandor en las aristas (solo para las delanteras)
\begin{scope}[on background layer]
  % Resplandor para aristas verticales delanteras
  \foreach \v in {A, Bp, D}{
    \draw[neonpink!30, line width=6pt, opacity=0.3] (T) -- (\v);
    \draw[neonpink!30, line width=6pt, opacity=0.3] (B) -- (\v);
  }
  
  % Resplandor para aristas horizontales delanteras del cuadrado
  \foreach \A/\B in {A/Bp, Bp/C}{
    \draw[neonyellow!30, line width=5pt, opacity=0.3] (\A) -- (\B);
  }
\end{scope}
\end{tikzpicture}
\end{center}
\end{document}

